\chapter{Codigo relevante}
\label{appx:codigo}

\section{Clase Subestacion}

\begin{lstlisting}[language=Python]
from dataclasses import dataclass, field
from typing import List

@dataclass
class Subestacion:
    S_nom_MVA: float = 30.0
    V_at_kV: float = 110.0
    V_bt_kV: float = 23.0
    X_pu: float = 0.12
    R_pu: float = 0.005
    P_fe_kW: float = 35.0
    P_cu_kW_nom: float = 180.0
    tau_termico_min: float = 90.0

    n_lineas: int = 4
    R_linea_pu: List[float] = field(default_factory=lambda: [0.04, 0.06, 0.05, 0.07])
    X_linea_pu: List[float] = field(default_factory=lambda: [0.08, 0.10, 0.09, 0.11])

    V_min_pu: float = 0.93
    V_max_pu: float = 1.07
    I_max_A: float = 800.0
    f_min_Hz: float = 49.5
    f_max_Hz: float = 50.5
    T_max_dev_C: float = 95.0

    P_nom_MW: List[float] = field(default_factory=lambda: [6.0, 4.5, 5.0, 3.5])
    Q_nom_MVAR: List[float] = field(default_factory=lambda: [1.5, 1.2, 1.3, 0.9])
\end{lstlisting}

\section{Clase GemeloDigital}

\begin{lstlisting}[language=Python]
class GemeloDigital:
    def __init__(self, se: Subestacion, dt_s: float = 1.0):
        self.se = se
        self.dt = dt_s
        self.V_at = 1.0
        self.V_bt = 1.0
        self.f = 50.0
        self.omega = 1.0
        self.T_dev = 25.0
        self.tap = 0.0

    def paso_estado_estacionario(
        self, P_carga_MW: np.ndarray,
        Q_carga_MVAR: np.ndarray,
        V_ref_pu: float = 1.0,
    ) -> Dict[str, float]:
        V_bt = np.clip(V_ref_pu + 0.02 * self.tap,
                        self.se.V_min_pu, self.se.V_max_pu)
        I = np.zeros(self.se.n_lineas)
        for k in range(self.se.n_lineas):
            V_k = V_bt - (self.se.R_linea_pu[k] * P_carga_MW[k]
                          + self.se.X_linea_pu[k] * Q_carga_MVAR[k]) \
                       / max(self.se.S_nom_MVA, 1e-3)
            V_k = max(V_k, 0.7)
            S_k = math.sqrt(P_carga_MW[k]**2 + Q_carga_MVAR[k]**2)
            I[k] = S_k * 1e3 / (math.sqrt(3) * self.se.V_bt_kV * V_k)
        P_total = P_carga_MW.sum()
        Q_total = Q_carga_MVAR.sum()
        S_total = math.sqrt(P_total**2 + Q_total**2)
        P_cu = 2 * self.se.P_cu_kW_nom * (S_total / (2 * self.se.S_nom_MVA))**2 / 1e3
        P_fe = 2 * self.se.P_fe_kW / 1e3
        P_at = P_total + P_cu + P_fe
        f = 50.0 - 0.005 * max(0.0, P_at - 25.0)
        f = float(np.clip(f, 49.0, 51.0))
        return {
            'V_at_pu': 1.0,
            'V_bt_pu': float(V_bt),
            'f_Hz': f,
            'I_linea_A': I,
            'P_total_MW': float(P_total),
            'Q_total_MVAR': float(Q_total),
            'P_cu_MW': float(P_cu),
            'P_fe_MW': float(P_fe),
            'S_total_MVA': float(S_total),
        }
\end{lstlisting}

\section{Clase LSTMAutoencoder}

\begin{lstlisting}[language=Python]
import torch
import torch.nn as nn

class LSTMAutoencoder(nn.Module):
    def __init__(self, n_features: int, latent_dim: int = 16,
                 hidden: int = 64, window: int = 30):
        super().__init__()
        self.window = window
        self.encoder = nn.LSTM(n_features, hidden, num_layers=2,
                                batch_first=True, dropout=0.1)
        self.bottleneck = nn.Linear(hidden, latent_dim)
        self.decoder = nn.LSTM(latent_dim, hidden, num_layers=2,
                                batch_first=True, dropout=0.1)
        self.out = nn.Linear(hidden, n_features)

    def forward(self, x):
        h, _ = self.encoder(x)
        z = self.bottleneck(h[:, -1, :])
        z_rep = z.unsqueeze(1).repeat(1, self.window, 1)
        d, _ = self.decoder(z_rep)
        return self.out(d)
\end{lstlisting}

\section{Conformal prediction}

\begin{lstlisting}[language=Python]
ALPHA = 0.05

err_train = ((Xw_n_train - model(Xw_n_train).numpy()) ** 2).mean(axis=(1, 2))
q_hat = np.quantile(err_train, 1 - ALPHA)

def predecir_alerta(sample):
    err = ((sample - model(sample).numpy()) ** 2).mean()
    return err > q_hat
\end{lstlisting}

\section{FedAvg}

\begin{lstlisting}[language=Python]
def fedavg(states, pesos):
    total = sum(pesos)
    avg = {k: torch.zeros_like(v) for k, v in states[0].items()}
    for s, p in zip(states, pesos):
        for k in avg:
            avg[k] += s[k] * (p / total)
    return avg
\end{lstlisting}

\section{Catalogo de fallas}

\begin{lstlisting}[language=Python]
CATALOGO_FALLAS = {
    'F1_corto_mono':  {'severidad': 0.85, 'duracion_s': (0.1, 2.0)},
    'F2_sobrecarga':   {'severidad': 0.45, 'duracion_s': (60, 600)},
    'F3_sobret_AT':    {'severidad': 0.90, 'duracion_s': (0.5, 5.0)},
    'F4_degrad_aisl':  {'severidad': 0.30, 'duracion_s': (1800, 7200)},
    'F5_perd_com_PMU': {'severidad': 0.50, 'duracion_s': (10, 600)},
    'F6_desbalance':   {'severidad': 0.35, 'duracion_s': (300, 3600)},
    'F7_falla_buje':   {'severidad': 0.25, 'duracion_s': (3600, 14400)},
    'F8_evento_clima': {'severidad': 0.70, 'duracion_s': (600, 14400)},
}
\end{lstlisting}

El codigo completo (notebook ejecutable y dashboard) se entrega como
material complementario en el repositorio del proyecto.